\subsection{STRUCTURE OF THE COGEBRA U(g) IN CHARACTERISTIC 0}

In this no., K will be assumed to be a field of characteristic 0.

Let \(S(\mathfrak{g})\) be the symmetric algebra of the vector space \(\mathfrak{g}\) , \(c_{\mathrm{s}}\) its coproduct (Algebra, Chapter III, § 11, no. 1, Example 6) and \(\eta\) the canonical isomorphism of the vector space \(\mathsf{S}(\mathfrak{g})\) onto the vector space U (Chapter I, § 2, no. 7). Recall that if \(x_{1},\ldots ,x_{n}\) are in \(\mathfrak{g}\) , then

\[
\eta (x _ {1} \dots x _ {n}) = \frac {1}{n !} \sum_ {\tau \in \mathfrak {S} _ {n}} \sigma (x _ {\tau (1)}) \dots \sigma (x _ {\tau (n)}).\tag{9}
\]

In particular, for \(x \in g\) and \(n \geqslant 0\) ,

\[
\eta (x ^ {n}) = \sigma (x) ^ {n}.\tag{10}
\]

Note that by Algebra, Chapter III, § 6, no. 1, Remark 3, \(\eta\) is the unique linear mapping of \(\mathbb{S}(\mathfrak{g})\) into U satisfying condition (10).

PROPOSITION 9. For every integer \(n \geqslant 0\) , let \(\mathbf{U}^n\) be the vector subspace of \(\mathbf{U}\) generated by the \(\sigma(x)^n\) for \(x \in \mathfrak{g}\) .

\begin{enumerate}
\def\labelenumi{(\alph{enumi})}
\tightlist
\item
  The sequence \((\mathbf{U}^n)_{n\geqslant 0}\) is a graduation of the vector space \(\mathbf{U}\) compatible with its cogebra structure.
\end{enumerate}

Let U be given the graduation (U \(^{n}\) ).

\begin{enumerate}
\def\labelenumi{(\alph{enumi})}
\setcounter{enumi}{1}
\tightlist
\item
  The canonical mapping \(\eta: \mathbf{S}(\mathfrak{g}) \to \mathbf{U}\) is an isomorphism of graded cogebras. Let \(x \in \mathfrak{g}\) and \(n \in \mathbf{N}\) . Then
\end{enumerate}

\[
c _ {\mathrm{S}} (x ^ {n}) = c _ {\mathrm{S}} (x) ^ {n} = (x \otimes 1 + 1 \otimes x) ^ {n} = \sum_ {i = 0} ^ {n} {\binom {n} {i}} x ^ {i} \otimes x ^ {n - i}\tag{11}
\]

since \(c_{\mathbb{S}}\) is an algebra homomorphism. Similarly, by (10),

\[
\begin{array}{l} c (\eta (x ^ {n})) = c (\sigma (x) ^ {n}) = c (\sigma (x)) ^ {n} = (\sigma (x) \otimes 1 + 1 \otimes \sigma (x)) ^ {n} \\ = \sum_ {i = 0} ^ {n} \binom {n} {i} \sigma (x) ^ {i} \otimes \sigma (x) ^ {n - i} = \sum_ {i = 0} ^ {n} \binom {n} {i} \eta (x ^ {i}) \otimes \eta (x ^ {n - i}), \end{array}\tag{12}
\]

whence

\[
(\eta \otimes \eta) (c _ {\mathrm{S}} (x ^ {n})) = c (\eta (x ^ {n})).
\]

As the \(x^n\) , for \(x \in \mathfrak{g}\) and \(n \in \mathbf{N}\) , generate the vector space \(\mathsf{S}(\mathfrak{g})\) , \((\eta \otimes \eta) \circ c_{\mathbf{S}} = c \circ \eta\) and \(\eta\) is a cogebra isomorphism.

On the other hand, formula (10) shows that \(\eta(\mathbb{S}^n(\mathfrak{g})) = \mathrm{U}^n\) , which completes the proof of (a) and (b) taking account of the fact that the graduation of \(\mathbb{S}(\mathfrak{g})\) is compatible with its cogebra structure.

The graduation \((\mathrm{U}^{n})_{n\geqslant0}\) of U is called the canonical graduation.

COROLLARY. The canonical mapping \(\sigma\) defines an isomorphism of \(g\) onto the Lie algebra \(\mathrm{P}(\mathbf{U})\) of primitive elements of \(\mathbf{U}\) .

As \(c^{+}\) is a graded homomorphism of degree 0,

\[
\mathrm{P} (\mathrm{U}) = \sum_ {n \geqslant 1} (\mathrm{P} (\mathrm{U}) \cap \mathrm{U} ^ {n}).
\]

It suffices to prove that if \(n > 1\) and \(a \in \mathrm{U}^n\) is primitive, then \(a = 0\) . Now \(a\) can be written as \(\sum_{i} \lambda_i a_i^n\) , where \(\lambda_i \in \mathbf{K}\) , \(a_i \in \sigma(g)\) . By (12), the term of bi-degree \((1, n - 1)\) in \(c^+(a)\) is \(n \sum_{i} \lambda_i a_i \otimes a_i^{n-1}\) . Hence \(\sum_{i} \lambda_i a_i \otimes a_i^{n-1} = 0\) . If \(\mu: \mathrm{U} \otimes \mathrm{U} \to \mathrm{U}\) is the linear mapping defined by multiplication on \(\mathrm{U}\) , then

\[
a = \sum_ {i} \lambda_ {i} a _ {i} ^ {n} = \mu \left(\sum_ {i} \lambda_ {i} a _ {i} \otimes a _ {i} ^ {n - 1}\right) = 0.
\]

Remarks. (1) \(\mathrm{U}_n = \sum_{i=0}^{n} \mathrm{U}^i\) (Chapter I, § 2, no. 7, Corollary 4 to Theorem 1).

\begin{enumerate}
\def\labelenumi{(\arabic{enumi})}
\setcounter{enumi}{1}
\tightlist
\item
  The mapping \(\eta\) is the unique morphism of graded cogebras of \(\mathsf{S}(\mathfrak{g})\) into U such that \(\eta(1) = 1\) and \(\eta(x) = \sigma(x)\) for \(x \in \mathfrak{g}\) . For if \(\eta'\) is a morphism satisfying these conditions, we prove by induction on \(n\) that \(\eta'(x^n) = \eta(x^n)\) for \(x \in \mathfrak{g}\) and \(n > 1\) . As \(c_{\mathbb{S}}^{+}(x^{n}) = \sum_{i=1}^{n-1} \binom{n}{i} x^{i} \otimes x^{n-i}\) by (3) and (11),
\end{enumerate}

\[
(\eta \otimes \eta) (c _ {\mathbf {S}} ^ {+} (x ^ {n})) = (\eta^ {\prime} \otimes \eta^ {\prime}) (c _ {\mathbf {S}} ^ {+} (x ^ {n}))
\]

by the induction hypothesis. It follows that \(c^{+}(\eta(x^{n})) = c^{+}(\eta'(x^{n}))\) ; it follows that \(\eta(x^{n}) - \eta'(x^{n})\) is a primitive element of degree \(n\) and hence is zero (Corollary to Proposition 9).

\begin{enumerate}
\def\labelenumi{(\arabic{enumi})}
\setcounter{enumi}{2}
\tightlist
\item
  Let \(\psi\) be the canonical isomorphism of the bigebra \(\mathsf{TS}(\mathfrak{g})\) onto the bigebra \(\mathsf{S}(\mathfrak{g})\) (Algebra, Chapter IV, § 5, Corollary 1 to Proposition 12). The mapping
\end{enumerate}

\[
\eta \circ \psi : T S (\mathfrak {g}) \rightarrow U
\]

is called canonical. It is the unique morphism \(\eta'\) of graded cogebras of TS(g) into U such that \(\eta'(1) = 1\) and \(\eta'(x) = \sigma(x)\) for all \(x \in \mathfrak{g}\) .

\begin{enumerate}
\def\labelenumi{(\arabic{enumi})}
\setcounter{enumi}{3}
\tightlist
\item
  Let V be a vector space. The primitive elements of the bigebra \(S(V)\) are the elements of degree 1. This follows from the Corollary to Proposition 9 applied to the commutative Lie algebra V.
\end{enumerate}

Let \((e_i)_{i\in I}\) be a basis of the vector K-space \(g\) , where the indexing set I is totally ordered. For all \(\alpha \in \mathbf{N}^{(I)}\) , we write

\[
e _ {\alpha} = \prod_ {i \in I} \frac {\sigma (e _ {i}) ^ {\alpha (i)}}{\alpha (i) !}.\tag{13}
\]

The \(e_{\alpha}\) , for \(|\alpha| \leqslant n\) , form a basis of the vector K-space \(\mathbf{U}_n\) (Chapter I, §2, no. 7, Corollary 3 to Theorem 1). Then

\[
e _ {0} = 1, \quad e _ {\varepsilon_ {i}} = \sigma (e _ {i}) \quad \text { for } i \in I.
\]

As the graded algebra associated with the filtered algebra U is commutative (loc. cit., Theorem 1), for \(\alpha, \beta\) in \(\mathbf{N}^{(I)}\) ,

\[
e _ {\alpha}. e _ {\beta} \equiv ((\alpha , \beta)). e _ {\alpha + \beta} \mod U _ {| \alpha | + | \beta | - 1}.\tag{14}
\]

where

\[
((\alpha , \beta)) = \prod_ {i \in I} \frac {(\alpha (i) + \beta (i)) !}{\alpha (i) ! \beta (i) !}.
\]

On the other hand, we have immediately

\[
\varepsilon (e _ {0}) = 1, \quad \varepsilon (e _ {\alpha}) = 0 \quad \text { for } | \alpha | \geqslant 1.\tag{15}
\]

Finally, formula (12) implies that, for \(\alpha \in \mathbf{N}^{(I)}\) ,

\[
c (e _ {\alpha}) = \sum_ {\beta + \gamma = \alpha} e _ {\beta} \otimes e _ {\gamma}.\tag{16}
\]

This formula allows us to determine the algebra \(\mathrm{U}' = \mathrm{Hom}(\mathrm{U},\mathrm{K})\) dual to the cogebra U (Algebra, Chapter III, § 11, no. 2). For let \(\mathrm{K}[[\mathrm{X}_i]_{i\in \mathbf{I}}]\) be the algebra of formal power series in indeterminates \((\mathrm{X}_i)_{i\in \mathbf{I}}\) (cf.~Algebra, Chapter III, § 2, no. 11); if \(\lambda \in \mathrm{U}'\) , let \(f_{\lambda}\) denote the formal power series

\[
f _ {\lambda} = \sum_ {\alpha} \langle \lambda , e _ {\alpha} \rangle \mathrm{X} ^ {\alpha}, \text { where } \mathrm{X} ^ {\alpha} = \prod_ {i \in I} \mathrm{X} _ {i} ^ {\alpha (i)}
\]

and the summation index \(\alpha\) runs through \(\mathbf{N}^{(I)}\) .

PROPOSITION 10. The mapping \(\lambda \mapsto f_{\lambda}\) is an isomorphism of the algebra \(\mathbf{U}'\) onto the algebra of formal power series \(\mathbf{K}[[\mathbf{X}_i]]_{i \in I}\) .

Because \((e_{\alpha})\) is a basis of U, the mapping \(\lambda \mapsto f_{\lambda}\) is K-linear and bijective. On the other hand, for \(\lambda, \mu\) in U',

\[
\begin{array}{r l} f _ {\lambda \mu} & = \sum_ {\alpha} \langle \lambda \mu , e _ {\alpha} \rangle \mathbf {X} ^ {\alpha} = \sum_ {\alpha} \langle \lambda \otimes \mu , c (e _ {\alpha}) \rangle \mathbf {X} ^ {\alpha} \\ & = \sum_ {\alpha} \langle \lambda \otimes \mu , \sum_ {\beta + \gamma = \alpha} e _ {\beta} \otimes e _ {\gamma} \rangle \mathbf {X} ^ {\alpha} \\ & = \sum_ {\beta , \gamma} \langle \lambda , e _ {\beta} \rangle \langle \mu , e _ {\gamma} \rangle \mathbf {X} ^ {\beta + \gamma} = f _ {\lambda} f _ {\mu}, \end{array}\tag{by (16)}
\]

which shows that \(\lambda \mapsto f_{\lambda}\) is an algebra isomorphism and completes the proof.

